% Algebraic Genomics -- the master volume.
% Build (from monograph/):  pdflatex AG_monograph; bibtex AG_monograph; pdflatex AG_monograph (x2)
% Chapters are read from chapters/chNN.tex, produced from papers/ by make_chapters.py (step 3 of the
% roadmap).  Until a chapter has been converted, a placeholder page is printed in its place.
\documentclass[11pt,twoside,openright]{book}
\usepackage{lmodern}
\usepackage[T1]{fontenc}
\usepackage[utf8]{inputenc}
\usepackage{amsmath,amssymb,amsthm,mathtools}
\usepackage[margin=1.05in]{geometry}
\usepackage{booktabs,array,longtable,multirow}
\usepackage{graphicx}
\usepackage{xcolor}
\usepackage{tikz}
\usetikzlibrary{arrows.meta,positioning}
\usepackage{microtype}
\usepackage{url}
\usepackage[hidelinks]{hyperref}
\usepackage{fancyhdr}
\raggedbottom
\emergencystretch=2em
\setlength{\headheight}{14pt}
\pagestyle{fancy}
\fancyhf{}
\fancyhead[LE]{\small\slshape\nouppercase{\leftmark}}
\fancyhead[RO]{\small\slshape\nouppercase{\rightmark}}
\fancyfoot[C]{\thepage}
\renewcommand{\headrulewidth}{0pt}
\fancypagestyle{plain}{\fancyhf{}\fancyfoot[C]{\thepage}\renewcommand{\headrulewidth}{0pt}}
\setcounter{tocdepth}{1}
\setcounter{secnumdepth}{2}

% ---- numbering: theorems, equations and figures by chapter
\numberwithin{equation}{chapter}
\theoremstyle{plain}
\newtheorem{theorem}{Theorem}[chapter]
\newtheorem{proposition}[theorem]{Proposition}
\newtheorem{lemma}[theorem]{Lemma}
\newtheorem{corollary}[theorem]{Corollary}
\newtheorem{conjecture}[theorem]{Conjecture}
\theoremstyle{definition}
\newtheorem{definition}[theorem]{Definition}
\newtheorem{example}[theorem]{Example}
\theoremstyle{remark}
\newtheorem{remark}[theorem]{Remark}
\newenvironment{chaptersummary}{\begin{quotation}\small\noindent\textbf{Summary.}\ }{\end{quotation}\medskip}

% ---- release metadata
\newcommand{\AGdoi}{10.5281/zenodo.22943487}
\newcommand{\AGrelease}{24 September 2026}
\newcommand{\AGversion}{1.0.0}

% ---- unified notation (AG_macros.tex collects every chapter's macros, conflicts resolved)
\InputIfFileExists{AG_macros.tex}{}{%
  \newcommand{\Z}{\mathbb{Z}}\newcommand{\Q}{\mathbb{Q}}\newcommand{\K}{\mathcal K}
  \newcommand{\DS}{\mathrm{DS}}\newcommand{\Lm}{\Lambda^{-}}\newcommand{\Nseq}{N_{\mathrm{seq}}}
  \newcommand{\ec}{\operatorname{ec}}\newcommand{\Loc}{\operatorname{Loc}}
  \newcommand{\rc}{\operatorname{rc}}\newcommand{\spec}{\operatorname{spec}}
  \newcommand{\sP}{\#\mathrm{P}}}

% ---- clickable DOIs in the bibliography
\providecommand{\doi}[1]{\href{https://doi.org/#1}{\nolinkurl{doi:#1}}}

% ---- a chapter, or a placeholder until it has been converted
\newcommand{\AGchapter}[3]{% #1 file, #2 title, #3 source paper
  \IfFileExists{chapters/#1.tex}{\input{chapters/#1}}{%
    \chapter{#2}\label{chap:#1}
    \sloppy\noindent\emph{Placeholder.} This chapter is the paper #3, to be converted by
    \texttt{make\_chapters.py}; the standalone paper is in \texttt{papers/}.}}

\begin{document}
\frontmatter
\begin{titlepage}
\centering
\vspace*{3cm}
{\Huge\bfseries Algebraic Genomics\par}
\vspace{0.6cm}
{\Large The Combinatorics of Sequence Reconstruction\\[2pt] and Representation\par}
\vspace{2cm}
{\Large Ruqing Chen\par}
\vfill
{\large Exact Eulerian reconstruction, double-stranded spectral fibres,\\
sandpile groups, and the computational complexity of sequence space\par}
\vspace{1.5cm}
{\small
\begin{tabular}{rl}
DOI: & \href{https://doi.org/\AGdoi}{\texttt{\AGdoi}}\\[2pt]
Archive: & Zenodo, \href{https://doi.org/\AGdoi}{\texttt{https://doi.org/\AGdoi}}\\[2pt]
Code and data: & \href{https://github.com/Ruqing1963/algebraic-genomics}{\texttt{github.com/Ruqing1963/algebraic-genomics}}\\[2pt]
Release: & version \AGversion, \AGrelease
\end{tabular}\par}
\vspace{1.2cm}

{\small Texts and figures: CC BY 4.0.\quad Code: MIT License.\par}
\end{titlepage}
\tableofcontents
% Preface of the monograph (input by AG_monograph.tex)
\chapter*{Preface}
\addcontentsline{toc}{chapter}{Preface}
\markboth{Preface}{Preface}

\paragraph{What this book is about.} Sequencing does not read a genome; it reads its words. From
short reads one learns, at best, the multiset of $k$-mers of the genome and of its reverse
complement, and everything an assembler, a variant caller or a sequence encoder does is constrained
by what that multiset determines. This book is about the exact size and shape of that constraint.
Given the data, the set of genomes compatible with them is a \emph{fibre}. We count fibres exactly,
describe the moves that connect their members, prove where exact counting stops being tractable, and
show where the same fibres reappear in modern machine-learning representations of biological
sequences.

\paragraph{One strand (Part I).} For a single strand the fibre is the set of Eulerian circuits of a
de Bruijn multigraph, and its size is a determinant, by the theorem of van Aardenne-Ehrenfest,
de Bruijn, Smith and Tutte. The determinant factors into a local part, fixed by the branch
statistics, and a global part: the order of the graph's sandpile group. Chapter~1 shows that the
global part is not determined by the branch statistics. Chapter~2 turns the group into a coordinate
system on reconstructions, one chip per rearrangement. Chapter~3 splits it along the repeats into a
part that longer words remove and a part that no word length removes. Chapter~4 shows where pairwise
repeat data stop describing it.

\paragraph{Two strands (Part II).} Real data do not know which strand a word came from. The
reverse-complement double cover keeps every single-strand lemma (Chapter~5). But its involution is an
anti-automorphism of the Laplacian (Chapter~6), and the double-stranded count is no longer one
determinant but a sum of determinants over the lattice points of a box (Chapter~7). For the phage
$\phi$X174 at word length $10$ the box has $2\,481\,687\,900$ points (Chapter~8), and the count is
$31\,925\,753\,246\,212$ (Chapter~9). Chapter~10 proves that this is not a failure of ingenuity:
the double-stranded count is $\#$P-hard, so no closed formula computes it unless $\mathrm{FP}=\#\mathrm P$.

\paragraph{Fibres and moves (Part III).} The members of a fibre are connected by explicit moves:
transpositions between direct repeats and inversions between inverted repeats. Chapter~11 computes
the lattice that the inversions span, by a Tate-cohomology dichotomy, and states a conjecture on
generation. Chapter~12, written with Zhengyi Chen, shows that the fibres are exactly what local
sequence encoders cannot see. It computes them for the whole human proteome and constructs explicit
proteins that a widely used composition encoder cannot tell apart. Chapters~13 and~14 treat pooled
data: polyploid phasing as a decomposition of one flow, and pan-genomes through their recombination
group.

\paragraph{Algorithms (Part IV).} Hardness is a statement about the worst case. Real genomes are
tractable when their tangled parts are small, and the book names the reason each time. Chapter~15
reduces the $4.6$-million-vertex graph of \emph{E.~coli} to $115$ vertices by contracting everything
except repeats. Chapter~16 computes the single-strand determinant in linear time along the snarl tree
of a genome graph. Chapter~17 describes the polytope of isoform abundances compatible with splicing
data.

\paragraph{How the results are checked.} Every chapter comes with a verification script, and every
number in the text, the tables and the figures is printed by one of them. The scripts check three
kinds of statement, and the book keeps them apart:
\begin{itemize}
\item \emph{theorems}, which are proved in the text; the scripts test them on exhaustive or random
instances, and a failure would indicate an error in the proof or in the code;
\item \emph{computations}, such as $\#\DS(\phi\mathrm X174,10)$, which the scripts produce and
cross-check by independent methods;
\item \emph{conjectures and empirical statements}, which are labelled as such and for which the
scripts report the evidence and nothing more.
\end{itemize}
The repository \url{https://github.com/Ruqing1963/algebraic-genomics}, archived under
DOI~\href{https://doi.org/\AGdoi}{\texttt{\AGdoi}}, contains the papers, the scripts, the input data
and the reference transcripts. The script \texttt{build\_all.py} compiles every chapter and reruns
every check; \texttt{-{}-quick} takes about half an hour, \texttt{-{}-full} an afternoon. The state of
that build for this release is recorded at the end of the book.

\paragraph{What the book corrects.} Several chapters began as specifications that turned out to be
wrong in part: a rank that was not the number of inverted repeats, a normalisation that was missing,
a tensor network that was a Schur complement, a polynomial formula that cannot exist. Each chapter
says which parts of its specification failed and why. We have kept these corrections in the text
because they are part of the result.

\paragraph{How to read it.} Part~I is self-contained. Part~II assumes Chapters~1 and~3. Part~III can be
read after Chapter~7. Part~IV uses Chapter~3 (Chapter~15) and the BEST theorem (Chapters~16, 17). The
chapters were written as papers; we have unified the notation, numbering and bibliography, in Part~I replaced repeated definitions by references to Chapter~1, and each
chapter still opens with a summary that can be read alone.

\paragraph{Acknowledgements.} Chapter~12 is joint work with Zhengyi Chen (Guangxi Key Laboratory of
Drug Discovery and Optimization, School of Pharmacy, Guilin Medical University), and is included with
the co-author's agreement. The phage, bacterial and human data are those of NCBI RefSeq, GenBank and UniProtKB.

\bigskip
\noindent Ruqing Chen\hfill Rivi\`ere-Beaudette, Qu\'ebec, \AGrelease


\mainmatter
\part{Classical foundations: the BEST theorem and the sandpile group}
\AGchapter{ch01}{Critical groups of de Bruijn graphs}{Bio~1}
\AGchapter{ch02}{The sandpile group as a coordinate system on reconstructions}{Bio~2}
\AGchapter{ch03}{The sandpile group splits along the repeats}{Bio~3}
\AGchapter{ch04}{Multi-copy repeats}{Bio~4}

\part{The double-stranded frontier and its complexity}
\AGchapter{ch05}{The reverse-complement double cover}{Bio~6}
\AGchapter{ch06}{The reverse-complement quotient is a signed graph}{Bio~7}
\AGchapter{ch07}{Double-stranded assembly as a lattice sum}{Bio~8}
\AGchapter{ch08}{Frontier elimination}{Bio~12}
\AGchapter{ch09}{The lattice sum as a fermionic partition function}{Bio~13}
\AGchapter{ch10}{Double-stranded reconstruction is \#P-hard}{Bio~19}

\part{Spectral fibres and rearrangements}
\AGchapter{ch11}{Inversions at inverted repeats}{Bio~17}
\AGchapter{ch12}{What a local encoder cannot see (with Zhengyi Chen)}{Chen \& Chen}
\AGchapter{ch13}{Polyploid phasing}{Bio~9}
\AGchapter{ch14}{The recombination group of a pan-genome}{Bio~10}

\part{Algorithmic implementations on genomes}
\AGchapter{ch15}{The bundle-chain decomposition of \emph{E.~coli}}{Bio~5}
\AGchapter{ch16}{The BEST determinant along the snarl tree}{Bio~14}
\AGchapter{ch17}{The isoform decomposition polytope}{Bio~11}

\backmatter
% Epilogue of the monograph (input by AG_monograph.tex)
\chapter*{Epilogue: what is open}
\addcontentsline{toc}{chapter}{Epilogue: what is open}
\markboth{Epilogue}{Epilogue}

Each problem below is left open by a specific chapter, and each is stated in the terms of that
chapter.

\paragraph{1. Generation of the double-stranded fibre.} Are any two double-stranded reconstructions
of a circular sequence joined by transpositions at repeated $k$-mers, inversions at inverted repeats
of length $k$, and global reverse complementation? This is true within one lattice point (Ukkonen,
Pevzner), true on every case that could be closed, and true on the palindromic class of Chapter~10 by
Kotzig's theorem (Chapter~11). The open part is the passage between lattice points: whether any two box
points with connected support are joined by \emph{realizable} inversions.

\paragraph{2. Odd word length.} Chapter~10's reduction needs palindromic $k$-mers, which exist only for
even $k$. For odd $k$ the quotient of the double cover is a signed graph without fixed vertices
(Chapter~6), and a reconstruction is an Eulerian circuit of a bidirected graph. Is the count still
$\#$P-hard?

\paragraph{3. Approximation.} Does the double-stranded count admit a fully polynomial randomized
approximation scheme? The same question for Eulerian circuits of undirected graphs is open, and
Chapter~10 shows that an answer for genomes would answer it for $4$-regular graphs.

\paragraph{4. Parity of flip distances.} In Chapter~7, reconstructions with the same spectrum are joined
by flip paths of even length only, at $k=11$ and $k=12$ for $\phi$X174. We have no explanation.

\paragraph{5. Beyond two copies.} The sandpile group of a repeat skeleton with two-copy repeats is
presented by an interlace matrix (Chapter~3). For three or more copies no pairwise presentation
exists (Chapter~4). Is there a presentation by data of higher order, and what is it?

\paragraph{What is not on this list.} The specifications of several chapters proposed further
directions that the chapters themselves do not support: Massey products on graphs, whose second
cohomology vanishes; a sheaf formalism for phasing, which leaves its complexity unchanged; and a gauge
theory of chromatin, which has no observable that can be derived. Thermodynamic phase transitions and
anyonic representations of sequences were also suggested, but no chapter defines the objects they would
be about. We list them here only to say that, in the form proposed, they are not open problems of this
book. A problem belongs on the list above only once it can be stated in the terms of a chapter.

% Guide to chapters and verification scripts (input by AG_monograph.tex, back matter)
\chapter*{Guide to chapters and verification scripts}
\addcontentsline{toc}{chapter}{Guide to chapters and verification scripts}
\markboth{Guide to chapters and verification scripts}{Guide to chapters and verification scripts}

Every chapter was first written as a paper with its own verification script. The scripts, their
reference transcripts and their input data live in the repository
\url{https://github.com/Ruqing1963/algebraic-genomics}, archived under
DOI~\href{https://doi.org/\AGdoi}{\texttt{\AGdoi}}, in the folder listed below. In the second
column, ``Bio~$n$'' is the number of the paper in its series; the numbers~15, 16 and~18 were
withdrawn before being written. \texttt{build\_all.py} recompiles every paper and reruns every
script: \texttt{-{}-quick} in about half an hour, \texttt{-{}-full} including the long computations.
The last column gives the checks of the reference transcript.

\medskip
{\small
\renewcommand{\arraystretch}{1.15}
\begin{longtable}{>{\raggedright\arraybackslash}p{0.33\textwidth}l>{\raggedright\arraybackslash}p{0.36\textwidth}r}
\toprule
chapter & paper & folder \texttt{papers/} and script & checks\\
\midrule
\endhead
\ref{chap:ch01}\ Critical groups of de Bruijn graphs & Bio~1 & \texttt{Ch01\_Bio\_01\_Critical\_Groups/} \texttt{dbg\_sandpile.py} & 30\\
\ref{chap:ch02}\ The sandpile group as a coordinate system & Bio~2 & \texttt{Ch02\_Bio\_02\_Rotor\_Routing/} \texttt{rotor\_assembly.py} & 8\\
\ref{chap:ch03}\ The sandpile group splits along the repeats & Bio~3 & \texttt{Ch03\_Bio\_03\_Repeat\_Splitting/} \texttt{repeat\_splitting.py} & 7\\
\ref{chap:ch04}\ Multi-copy repeats & Bio~4 & \texttt{Ch04\_Bio\_04\_Multicopy\_Repeats/} \texttt{multicopy.py} & 4\\
\ref{chap:ch05}\ The reverse-complement double cover & Bio~6 & \texttt{Ch05\_Bio\_06\_RC\_Double\_Cover/} \texttt{rc\_double\_cover.py} & 6\\
\ref{chap:ch06}\ The quotient is a signed graph & Bio~7 & \texttt{Ch06\_Bio\_07\_Signed\_Quotient/} \texttt{bidirected\_sandpile.py} & 9\\
\ref{chap:ch07}\ Double-stranded assembly as a lattice sum & Bio~8 & \texttt{Ch07\_Bio\_08\_DS\_Lattice\_Sum/} \texttt{ds\_lattice.py}, \texttt{ds\_flips.py} & 9, 4\\
\ref{chap:ch08}\ Frontier elimination & Bio~12 & \texttt{Ch08\_Bio\_12\_Frontier\_Elimination/} \texttt{ds\_transfer\_matrix.py} & 8\\
\ref{chap:ch09}\ A fermionic partition function & Bio~13 & \texttt{Ch09\_Bio\_13\_Fermionic\_Partition/} \texttt{ds\_fermion.py}, \texttt{ds\_count.py} & 7, 4\\
\ref{chap:ch10}\ Double-stranded reconstruction is \#P-hard & Bio~19 & \texttt{Ch10\_Bio\_19\_SharpP\_Hardness/} \texttt{bio19\_verify.py} & 5\\
\ref{chap:ch11}\ Inversions at inverted repeats & Bio~17 & \texttt{Ch11\_Bio\_17\_Inversions\_Fibre/} \texttt{bio17\_verify.py} & 6\\
\ref{chap:ch12}\ What a local encoder cannot see & joint & \texttt{Ch12\_Spectral\_Fibres/} \texttt{spectral\_fibres.py} & 5\\
\ref{chap:ch13}\ Polyploid phasing & Bio~9 & \texttt{Ch13\_Bio\_09\_Polyploid\_Phasing/} \texttt{polyploid\_phasing.py} & 6\\
\ref{chap:ch14}\ The recombination group of a pan-genome & Bio~10 & \texttt{Ch14\_Bio\_10\_Pangenome\_Sheaf/} \texttt{pangenome\_sheaf.py}, \texttt{real\_pangenome.py} & 5, 3\\
\ref{chap:ch15}\ The bundle-chain decomposition of \emph{E.~coli} & Bio~5 & \texttt{Ch15\_Bio\_05\_Ecoli\_Decomposition/} \texttt{ecoli\_sandpile.py} & 6\\
\ref{chap:ch16}\ The BEST determinant along the snarl tree & Bio~14 & \texttt{Ch16\_Bio\_14\_Snarl\_Schur/} \texttt{bio14\_verify.py} & 6\\
\ref{chap:ch17}\ The isoform decomposition polytope & Bio~11 & \texttt{Ch17\_Bio\_11\_Isoform\_Polytope/} \texttt{isoform\_polytope.py} & 7\\
\bottomrule
\end{longtable}}

\IfFileExists{AG_build_state.tex}{% Build record of this release (input by AG_monograph.tex, back matter)
\chapter*{Build record}
\addcontentsline{toc}{chapter}{Build record}
\markboth{Build record}{Build record}

The repository was built with \texttt{python build\_all.py -{}-quick} on 24 September 2026 (Python
3.12.4, Windows~11), after the chapter folders were renumbered. The build compiles every paper twice
with \texttt{pdflatex} and reruns its verification scripts; a script passes when it exits normally and
its final line reports all of its checks passed.

\begin{center}
\begin{tabular}{lr}
\toprule
items passed & 37\\
items failed & 0\\
items skipped (run only with \texttt{-{}-full}) & 4\\
wall-clock time & 24 min\\
\bottomrule
\end{tabular}
\end{center}

\noindent All $17$ papers compile with no error and no undefined reference. Every script run in quick mode
passes all of its checks. Two of them, Chapter~\ref{chap:ch08} (\texttt{ds\_transfer\_matrix.py}) and
Chapter~\ref{chap:ch09} (\texttt{ds\_count.py}), run a subset of their checks in quick mode ($7$ of
$8$, $2$ of $4$).

\paragraph{Not rerun in this build.} The four items that run only with \texttt{-{}-full} are the
signed-quotient battery of Chapter~\ref{chap:ch06}, the eight-strain pan-genome of
Chapter~\ref{chap:ch14}, and the proteome computation and its figures in Chapter~\ref{chap:ch12}. The
long computations of Chapters~\ref{chap:ch08} and~\ref{chap:ch09} were also run only in their quick
form. These are the connected count at $k=10$ and the enumeration of the $1\,240\,843\,950$ half-box
points. For all of them the numbers in the text come from the reference transcripts in the repository,
which were produced when the chapters were written. A \texttt{-{}-full} build, of about four hours,
reruns them.
}{}
\IfFileExists{AG_monograph.bib}{\bibliographystyle{plain}\bibliography{AG_monograph}}{}
\end{document}
